\scene{PREPARATION / NOT PART OF THE RECORDING}{A 19-minute Lea demonstration}{Leave one minute for compilation, network, or model latency.}

This script demonstrates the Lea extension inside Overleaf. Quoted paragraphs
are suggested narration; numbered steps are actions. Controls depend on the
item's current state. A reference proof or a finished example is not evidence
that an agent just completed a live run.

\cue{Files to keep open}
\begin{itemize}
  \item Upload \texttt{lea-overleaf-project.zip} as an Overleaf project. Keep
    \texttt{main.tex} as its main document and work in Code mode.
  \item Keep \texttt{copy-paste-snippets.txt} beside this script for the exact edits.
  \item Keep a second, independently rehearsed Overleaf project ready as a backup.
    A copied Overleaf project has a different Lea identity; prepare its Lea results
    separately rather than assuming history transfers.
\end{itemize}

\cue{Prepare the states before recording}
\begin{enumerate}
  \item Start the normal Lea services, load the extension, compile the document,
    open the floating \textbf{L}, and refresh the pane. Expect \textbf{9 marked items}.
    The unmarked Doubling zero lemma and the unmarked Shift identity are intentional.
  \item Formalize \texttt{demo\_double}, then \texttt{demo\_double\_add}, then
    \texttt{demo\_double\_two}. Inspect the actual names, namespaces, types, and
    imports. Rehearse snippets F--H so the corollary genuinely depends on the lemma.
    Restore explicit arguments and a checked corollary for the recording.
  \item Leave Doubling zero unmarked for the live insertion. Leave the two optional
    input examples unfinished for the batch demonstration.
  \item Enable \textbf{Pause formalization when Lea needs author input} before
    individually formalizing the flawed square claim. Preserve its real finding
    or checked disproof. If paused, confirm there is a usable saved checkpoint.
  \item Formalize \texttt{demo\_shifted\_square}, then \texttt{demo\_shift\_one},
    then \texttt{demo\_shift\_two}. Require an actual shared helper, its checked
    status, and both consumers' imports before describing the example as reuse.
  \item Keep a past batch report available. Rehearse ZIP export. Configure a
    demonstration GitHub remote beforehand only if you intend to show a real push
    or import. Keep saved credentials concealed throughout the recording.
\end{enumerate}

\cue{Time checkpoints}

0--3 min: open and mark.\quad
3--6: formalize and review.\quad
6--8: chat and edit.\quad
8--11: fix source.\quad
11--13: repair.\quad
13--15: batch.\quad
15--17: shared results and Blueprint.\quad
17--19: settings, export, and close.

\textbf{Latency rule:} spend at most about 20 seconds waiting without something
useful to explain. Then say, ``That run is still working. Here is the same workflow
on a result I prepared earlier.'' Switch only to a genuinely prepared state.

\clearpage
\scene{0:00--3:00}{Open Lea and mark a statement}{Start on page 1 of the compiled document, with main.tex open.}

\click
\begin{enumerate}
  \item Click the floating \textbf{L} button (\textbf{Lean project pane}). Resize
    the pane so both source and item details remain readable.
  \item Expand \texttt{main.tex} and its section. Point out the other two source
    files and their progress indicators. Use \textbf{Go to source} on Doubling a sum.
  \item In \texttt{main.tex}, find \texttt{DEMO LIVE INSERT} inside the Doubling zero
    lemma. Paste snippet A as one new line:
\end{enumerate}
\begin{lstlisting}
% lea: formalize label=demo_double_zero uses={demo_double}
\end{lstlisting}
\begin{enumerate}[resume]
  \item Let Overleaf save. Refresh the Lea pane if the item has not appeared.
    Expect \textbf{10 marked items}. Click the source badge and show
    \textbf{Show in Lean pane}, then return to the item's details.
  \item Point to the definition's \texttt{\string\leadefinition}, the custom
    \texttt{claim} environment's \texttt{\string\leatheorem}, and the corollary's
    \texttt{uses} and \texttt{context} fields. The package is already included in
    this project; explain \textbf{Add Lea tags} if that setup control is offered.
\end{enumerate}

\say
``This is a small mathematical document about doubling a natural number.
Lea adds a formalization workspace alongside the source. I choose which
mathematical blocks it should work on, and I can jump between the document
and the corresponding formal result.''

``This lemma was ordinary LaTeX. Adding one comment asks Lea to formalize it;
the comment does not change the PDF. Inline tags are another supported form.
They also let me use my own environment names, such as this Claim.''

``The dependency field points to earlier Lea items. The context field gives
guidance, such as notation or a proof strategy. Explicit labels make this demo
easy to follow; they are optional when I just want to mark a new item.''

\fallback
If the pane is empty, refresh after the source saves and check that the extension
is connected. Do not insert another marker in an already tagged block. If the live
marker is already present from rehearsal, remove that one line before recording.

\clearpage
\scene{3:00--6:00}{Formalize and inspect the evidence}{Live target: demo\_double\_zero. Prepared results: definition and additivity.}

\click
\begin{enumerate}
  \item Expand Doubling zero. Under \textbf{Other actions}, open the more-actions
    menu and choose \textbf{Stub} if offered. Show the skeleton and its unfinished
    obligation. This step is optional; \textbf{Formalize} works directly.
  \item Click \textbf{Formalize} (or \textbf{Re-formalize} after creating the stub).
    Expand the live progress/status view. Point out \textbf{Pause}; use it only if
    the run is still active and the pause/resume path was rehearsed.
  \item While it runs, expand the prepared Doubling definition and show that it
    produces a Lean definition with a body, not a theorem requiring a proof.
  \item Return to Doubling zero when it settles, or explicitly switch to the
    prepared Doubling a sum result. Show \textbf{Lean code} and its copy control.
  \item Inspect \textbf{Lean Check}. Expand \textbf{Lea Status details}, read the
    assessment's scope and any findings, then open \textbf{Update history}.
  \item Compare the Lean proposition with the LaTeX. On an eligible checked
    result you have actually reviewed, click the personal \textbf{checkmark}.
\end{enumerate}

\say
``A stub gives me a formal statement to work from, but it is still unfinished.
Formalize asks Lea to complete the task. I can follow its progress, and pause
an individual run when necessary.''

``Here is the resulting Lean. I check the types, hypotheses, and conclusion
against the mathematical statement. In this example the domain is the natural
numbers and the conclusion says that doubling zero gives zero.''

``These two indicators answer different questions. Lean Check reports the
formal checking result. Lea Status reports Lea's assessment of the source and
formalization, including findings or remaining obligations. I read the scope:
an assessment of partial work is not a completed proof.''

``This checkmark records my personal review of this revision and its
dependencies. It is separate from formal checking, and later changes can
invalidate it.''

\fallback
If a run finishes before Pause can be used, leave the completed result intact.
Show a genuine prepared paused item later. Do not call an unchecked artifact
verified, or describe confidence alone as a correctness guarantee.

\clearpage
\scene{6:00--8:00}{Ask a question, then edit the proof}{Prepared target: demo\_double\_add with explicit arguments m and n.}

\click
\begin{enumerate}
  \item Expand Doubling a sum. In \textbf{Other actions}, click \textbf{Chat}.
    Paste snippet B and send it:
\end{enumerate}
\begin{lstlisting}
Explain how the Lean statement matches the LaTeX claim and how the proof
rearranges addition. Do not change any files.
\end{lstlisting}
\begin{enumerate}[resume]
  \item Read a short portion of the response if it arrives. Point out
    \textbf{Open in LeaChat}, which continues the same session in the larger UI.
    Keep the recording focused on the extension.
  \item Wait until the item's chat run is idle before editing. Choose
    \textbf{Edit} from the more-actions menu. Preserve the generated imports,
    namespace, declaration name, and statement. Replace only the proof after
    \texttt{:=} with snippet C:
\end{enumerate}
\begin{lstlisting}
by
  simp only [demo_double]
  ring
\end{lstlisting}
\begin{enumerate}[resume]
  \item Click \textbf{Save}. Inspect the checking result and dependent impact
    summary. If the generated definition has a qualified name, use the exact name
    rehearsed for this project. Point out \textbf{Cancel} for discarding a draft.
\end{enumerate}

\say
``I can ask questions in the context of this item. Here I want an explanation
of how the formal statement and proof correspond to the source. The session
also remains available in LeaChat.''

``I can edit the Lean myself. This replacement keeps the same mathematical
statement: it unfolds doubling and lets the ring tactic rearrange the
arithmetic. Saving records and checks my edit, including affected dependents
when needed. It does not start another agent proof search.''

\fallback
If chat is slow, show the prepared conversation and identify it as earlier work.
Do not edit during its active run. If the generated names or statement differ
from the rehearsed version, use a prepared safe edit or show Edit/Cancel instead
of pasting a mismatched declaration.

\clearpage
\scene{8:00--11:00}{Find a source problem and correct it}{Target: demo\_square\_bound in sections/02-source-review.tex.}

\click
\begin{enumerate}
  \item Use \textbf{Go to source} on Square versus the original number. Show
    the original claim and its associated proof, marked as a flawed exercise.
  \item Open the genuine finding prepared before recording. If demonstrating a
    new run, first enable \textbf{Pause formalization when Lea needs author input}
    in the gear menu, save, and start this item individually.
  \item Expand \textbf{Lea Status details}. Explain the failed sign argument.
    If the outcome is \textbf{Claim disproved}, inspect whether the counterexample
    itself is checked. If paused, show the actual \textbf{Resume} control.
  \item In the source editor, replace the whole theorem and following proof
    with \textbf{snippet D}. This adds $x\ge1$ and corrects the sign argument.
    Keep the Lea label \texttt{demo\_square\_bound} unchanged. Let Overleaf save.
  \item For a compatible paused run, click \textbf{Resume}. For a completed
    disproof or another settled run, choose \textbf{Re-formalize}. These are
    different actions; describe the one actually used.
  \item Inspect the new Lean statement: it must contain the new hypothesis
    $x\ge1$. If a genuine review candidate is offered, inspect it before
    \textbf{Approve result}. That applies the reviewed candidate; it is separate
    from the personal review checkmark.
\end{enumerate}

\say
``The original statement says that every real number is at most its square.
That is false: at one half, the square is one quarter. The author's proof
also incorrectly assumes that both factors are always nonnegative.''

``Lea should surface the issue rather than quietly insert the missing
hypothesis. Depending on the run, that can appear as a request for author
input or as a checked counterexample to the original claim.''

``I am now changing the mathematics explicitly. With the assumption that
x is at least one, both factors are nonnegative and the proof works.
I check that the new formal statement includes that assumption.''

\cue{Optional, if time remains}
After the corrected result is checked, use snippet E to rewrite the displayed
conclusion as $x^2-x\ge0$. Observe outdated status, then \textbf{Re-formalize}.
Show that an earlier personal review or assessment does not stay current
merely because a proof was once checked.

\fallback
If Resume reports an incompatible checkpoint, keep the message visible and use
Re-formalize as a fresh attempt. This example has an author proof, so it is not
the example for the conditional \textbf{Continue best effort} action.

\clearpage
\scene{11:00--13:00}{Inspect a broken dependent and repair it}{Use the rehearsed additivity-to-corollary dependency.}

\click
\begin{enumerate}
  \item Confirm that Doubling a sum and Adding two before doubling are checked
    and idle. Open the corollary's Lean code. It must really apply
    \texttt{demo\_double\_add n 2}; snippet F prepares that proof if needed.
  \item Edit \texttt{demo\_double\_add}. Change only its binder notation from
    \texttt{(m n : Nat)} to \texttt{\{m n : Nat\}} as in snippet G. Preserve its
    proposition, name, proof body, imports, and namespace. Click \textbf{Save}.
  \item Inspect the impact summary and the dependent corollary's error. The
    lemma remains mathematically the same, but its old explicit application is
    incompatible with the new implicit arguments.
  \item Click \textbf{Repair with Lea} on the affected corollary. Point out
    \textbf{Repair all} if a group is offered. Inspect the resulting code and
    current Lean Check when the repair settles.
\end{enumerate}

\say
``These results are connected by actual Lean dependencies. I am making a
small interface change: the lemma's two arguments become implicit. Its
mathematical content stays the same, but an existing caller may need updating.''

``Lea reports the affected work instead of continuing to present its old
check as current. I can ask Lea to repair the dependent result and then
inspect what changed.''

\cue{Deterministic manual fallback}
If agent repair is slow, use a previously completed repair. Alternatively,
when no agent run is active, edit the dependent lemma application as in snippet H:
\begin{lstlisting}
-- Before:
demo_double_add n 2
-- After:
demo_double_add (m := n) (n := 2)
\end{lstlisting}
Click \textbf{Save} and inspect the check. Say, ``Here I repaired the call
manually; Repair with Lea automates that kind of adaptation.''

\fallback
This is an optional live breakage demonstration, not a required production
workflow. If the generated corollary does not call the lemma, do not claim the
edit broke it. If the save is rejected or the calling convention was not
rehearsed, show the prepared impact and repair example. Avoid leaving a live
repair active when starting the next batch.

\clearpage
\scene{13:00--15:00}{Run a batch and inspect its report}{Small remaining targets: standalone proposition and existing Lean.}

\click
\begin{enumerate}
  \item Open the optional input forms near the bottom of \texttt{main.tex}.
    Point out the standalone proposition tag and the \texttt{leacode} block.
    These are already marked, but can be left unfinished until this step.
  \item Return to the top of \textbf{Items}. Point out \textbf{Stub all} if
    eligible theorems remain. Click \textbf{Formalize all} when the project is
    idle and the source corrections and repairs have settled.
  \item Show the current item, queued work, and outcomes. Explain that eligible
    items are ordered by declared dependencies and completed work is not simply
    restarted. If shared consumers are still unfinished, expect their supporting
    work to appear as well.
  \item Point out \textbf{Stop} while the batch is active. Demonstrate it only
    if that stop/continuation was rehearsed. Use \textbf{Continue remaining} only
    when the actual paused batch offers it.
  \item When the batch settles, open \textbf{Formalize all report}. If it is
    still running, open \textbf{Past Formalize all reports} and identify the
    selected report as an earlier batch. Inspect outcomes and shared-result counts.
\end{enumerate}

\say
``For a larger document I can formalize eligible items as a batch, following
their declared dependencies. The panel shows what is running, what remains,
and what needs attention.''

``A finished batch is not a promise that every item succeeded. I read the
report for checked results, failures, findings, and shared mathematics.
Source findings stay visible during Formalize all; the batch does not use
the individual-run pause-on-author-input behavior.''

``These two extra inputs show that I can mark standalone prose or supply
existing Lean code. Code appearing in the PDF is only source material;
I still inspect Lea's checking result.''

\fallback
If no eligible items remain, show the saved report rather than deliberately
invalidating a checked result. Do not claim the initial marked-item count is a
count of newly completed proofs. A stub with an unfinished obligation remains
unfinished even if a stub-generation batch succeeds.

\clearpage
\scene{15:00--17:00}{Reuse mathematics and view the Blueprint}{Prepared shared example: shifted squares.}

\click
\begin{enumerate}
  \item Open \texttt{sections/03-shared-results.tex}. Show the definition
    $Q_a(x)=(x-a)^2$, the unmarked Shift identity $Q_a(a+t)=t^2$, and its two
    consumers at $t=1$ and $t=2$.
  \item Expand \textbf{Shared mathematics}. Select the helper produced or reused
    by the rehearsed run. Inspect its full signature, checked state, canonical
    module/import, attempts, and \textbf{Used by} consumers.
  \item Open both consumer artifacts. Confirm they import and use that same
    public helper. If a genuine dependency wait is present, use
    \textbf{View supporting result} and explain automatic continuation. Point
    out \textbf{Refresh status}, \textbf{Stop waiting}, or \textbf{Continue when ready}
    only when offered in that state.
  \item Switch from \textbf{Items} to \textbf{Blueprint}. Click
    \textbf{Generate from formalized theorems} if a starter graph is needed.
    Select a doubling node and inspect its kind, status, and dependencies.
\end{enumerate}

\say
``Both of these results use the same supporting identity. The identity is
source mathematics here; it is not an extra tagged document target. Lea can
create a shared formal result and reuse it for multiple consumers.''

``I verify reuse in the shared record and the generated imports. If the
helper is already checked, immediate reuse is expected. If it has to be proved,
the requesting item can wait and continue when the supporting work is ready.''

``Blueprint gives me a project view of the mathematical dependencies. I can
move from individual proofs to the structure of the formalization and inspect
which pieces are checked or still need work.''

\fallback
The source guidance requests shared-helper use, but agent choices are not a
deterministic test harness. If the live run proves each consumer independently,
say so; it has not demonstrated shared reuse. Use the prepared genuine shared
record. Do not invent a new helper attempt when an existing one was reused.

\clearpage
\scene{17:00--19:00}{Show controls, export, and finish}{Keep this final tour short. Reserve 19:00--20:00 for spillover.}

\click
\begin{enumerate}
  \item Click the gear for \textbf{Lea settings and usage}. Show usage and
    maximum spend, then model selection, maximum turns, source mirroring,
    and the individual-run author-input pause option. Keep configured API keys hidden.
  \item Show the project name and Lean namespace. Open \textbf{Edit project name}
    or \textbf{Edit name}, inspect the preview, then cancel. Avoid renaming the
    rehearsed project during this recording.
  \item Open \textbf{Share}, then \textbf{Download .zip}. Confirm the Lean project
    export is offered or downloaded. Explain that it contains the Lean project,
    distinct from the LaTeX project used as source.
  \item Show the configured remote and \textbf{Push to GitHub}. Execute a push
    only to the intended demo repository if it was prepared for the recording.
    Otherwise identify this as a tour of the control.
  \item Show \textbf{Add Lean files from GitHub}. With a prepared repository,
    inspect an actual import preview and its skipped conflicts before deciding
    whether to confirm. Imported work still needs the recorded checking process.
  \item Close the panel and return to page 1 with a checked result visible.
\end{enumerate}

\say
``I can choose the model, review usage, and control the run and spend limits
from the extension. Source mirroring makes surrounding LaTeX available as
context. The project also has a name and Lean namespace.''

``The formal work is portable: I can download the Lean project, push it to
a configured GitHub repository, or preview Lean files to add from GitHub.
For an import I inspect the preview and subsequent checks.''

``We started with mathematical prose, selected the statements to formalize,
and inspected the Lean and the source findings. We then refined a proof,
corrected a claim, handled dependencies, and exported the resulting project.''

\fallback
Do not spend recording time configuring credentials or hunting for a repository.
If there is no prepared remote, show the available actions and explain them
without claiming a push or import occurred. The ZIP export is the concrete
portable result to finish on.

\clearpage
\scene{REHEARSAL NOTES}{Cuts, feature coverage, and evidence}{Use this page before recording, not as a slide.}

\cue{For a 15-minute version}
Keep the same sequence. Use these checkpoints:
\begin{itemize}
  \item 0--2: open and mark; 2--5: live formalization and review.
  \item 5--6: chat or a safe manual edit; 6--9: the source correction.
  \item 9--10: show the prepared repair outcome; 10--12: batch and its report.
  \item 12--14: shared mathematics and Blueprint; 14--15: ZIP export and close.
\end{itemize}
Cut live dependency breakage, the optional second source revision, and the
GitHub/configuration tour first. Keep the check/status distinction and the
author's review of the mathematical statement.

\cue{What is covered}
\begin{itemize}
  \item Source navigation, file grouping/progress, inline badges, comments,
    definition/theorem/corollary tags, a custom environment, standalone prose,
    and an existing-Lean block.
  \item Labels, declared dependencies, guidance, associated author proofs,
    optional stubs, individual formalization, and pause/resume controls.
  \item Lean Check, live Lea Status and history, outdated source, a counterexample
    or source finding, source corrections and new formalizations, conditional
    result review, and personal review tied to a revision.
  \item Item chat, the LeaChat session link, manual edits/checks, dependent impact,
    individual/group repair, batch actions, saved reports, shared helper reuse,
    conditional dependency-wait controls, and Blueprint.
  \item Model/usage/limits, source context, project identity, ZIP export, and the
    GitHub push/import interfaces.
\end{itemize}

\cue{Conditional features}
Not every state is expected in one recording. \textbf{Approve result} needs a
real review candidate. \textbf{Continue best effort} needs an eligible source
pause without an associated author proof; the flawed-claim example deliberately
has a proof. \textbf{Continue remaining} needs a paused batch. Waiting controls
need a real shared dependency wait. Mention these only in their proper context
or use separately rehearsed examples.

\cue{What was validated when preparing this kit}
The LaTeX projects are compiled, the marker inventory and proof associations
are checked against the extension parser, and the reference Lean proofs are
checked locally using the repository's Lean/Mathlib toolchain. The reference
file also checks the false claim's counterexample and the repaired implicit
argument call. These checks do not certify an unrun agent-generated artifact,
guarantee model timing, or replace the live Overleaf rehearsal above.
